\vfill%
\subsection{Control overall width of table}
\tikz{\node[text width=\linewidth-6mm,draw=green!70,rounded corners=2pt,fill=black!10!green!10,inner sep=3mm] {%
Is there a simple way to set the periodic table to text width, column width, etc.?
};}%
\\ [6pt]Yes, there is. It can be done using the \textsf{\large\textbackslash resizebox} command provided by the graphicx package (and also by the graphics package). For example:
\\ [2pt]\tikz{\node[text width=\linewidth-8pt,inner xsep=4pt,align=left,fill=black!10,rounded corners=2pt] %
{\begin{minipage}{\linewidth}%
\begin{verbatim}
\resizebox{\linewidth}{!}{\pgfPT}
\end{verbatim}
\end{minipage}};}%
\\ or
\\ [2pt]\tikz{\node[text width=\linewidth-8pt,inner xsep=4pt,align=left,fill=black!10,rounded corners=2pt] %
{\begin{minipage}{\linewidth}%
\begin{verbatim}
\resizebox{\linewidth}{!}{\pgfPT[show title=false]}
\end{verbatim}
\end{minipage}};}%
\\ [2pt]will produce a Periodic Table with the width of the current \textsf{\large\textbackslash linewidth}, whatever is its value (the text width, the column width, the width of a minipage,\ldots), and with the proper scaling of its height.
\\ There is no need of loading the graphicx package since pgf-PeriodicTable loads the tikz package, which in turn loads the graphicx package.
\\ [2pt]\tikz{\path[draw=green!70,rounded corners=2pt,fill=black!10!green!10,rounded corners=2pt] (0,0) rectangle ++(\textwidth,-4.5pt);}%
\subsection{Compact Periodic Table}
\tikz{\node[text width=\linewidth-6mm,draw=green!70,rounded corners=2pt,fill=black!10!green!10,inner sep=3mm] {%
Is there a way to put groups 1 and 2 really next to group 13 to 18? That would make the whole thing more compact. I sometimes need just the representative elements for teaching purposes.
};}%
\\ [6pt]Although it is not common usage, it can be done:
\\ [2pt]\tikz{\node[text width=\linewidth-8pt,inner xsep=4pt,align=left,fill=black!10,rounded corners=2pt] %
{\begin{minipage}{\linewidth}%
\begin{verbatim}
\documentclass[border=10pt]{standalone}
\usepackage{pgf-PeriodicTable}
\usepgfPTlibrary{colorschemes}
\pgfPTGroupColors{example}{G1=blue!50!white,G2=green!90!white}
\pgfPTsetLanguage{de}
\begin{document}
% \pgfPTstyle[show title=false, back color scheme=example,show legend=false]
% \pgfPT[Z list = G1]\foreach \n in {2,13,14,15,16,17,18} {%
%                               \pgfPT[show period numbers=false,Z list = G\n]%
%                               }% make sure there are no spaces between \pgfPT
% or
% \pgfPTstyle[show title=false,show period numbers=false,back color scheme=example,
%                 show legend=false]
% \pgfPT[show period numbers,Z list = G1]\pgfPT[Z list = G2]\pgfPT[Z list = p]% make
%                                                                 sure there are no spaces between \pgfPT
% or
    \pgfPTstyle[show title=false,show period numbers=false,back color scheme=example,
                    show legend=false]
    \pgfPT[show period numbers,Z list = s]\pgfPT[Z list = p]
\end{document}
\end{verbatim}
\end{minipage}};}%
\newpage%
\begin{center}
\scalebox{.6}{\includegraphics{manualfiles/pgfPT_user_compactPT.pdf}}
\end{center}
\tikz{\path[draw=green!70,rounded corners=2pt,fill=black!10!green!10,rounded corners=2pt] (0,0) rectangle ++(\textwidth,-4.5pt);}%
\endinput
